\documentclass[a4paper,11pt]{article}
\usepackage[utf8]{inputenc}
\usepackage[T1]{fontenc}
\usepackage[french]{babel}
\usepackage[left=1cm,right=1cm,top=.5cm,bottom=0cm]{geometry}
\usepackage{enumitem}
\usepackage{hyperref}
\usepackage{xcolor}
\usepackage{titlesec}
\usepackage{fontawesome5}
\usepackage{lmodern}
\usepackage[normalem]{ulem}

% --- Couleurs CEA ---
\definecolor{primary}{RGB}{0, 82, 147}
\definecolor{secondary}{RGB}{80, 80, 80}

% --- Configuration des liens ---
\hypersetup{
    colorlinks=true,
    linkcolor=primary,
    filecolor=primary,
    urlcolor=primary,
}

% --- Formatage des titres ---
\titleformat{\section}{\large\bfseries\color{primary}\uppercase}{}{0em}{}[\titlerule]
\titlespacing{\section}{0pt}{8pt}{5pt}

% --- Commandes personnalisées ---
\newcommand{\resumeItem}[1]{\item\small{#1}}
\newcommand{\resumeSubheading}[4]{
  \vspace{-1pt}\item
    \begin{tabular*}{0.97\textwidth}[t]{l@{\extracolsep{\fill}}r}
      \textbf{#1} & \textit{\small #2} \\
      \textit{\small #3} & \textit{\small #4} \\
    \end{tabular*}\vspace{-5pt}
}

\begin{document}

% --- EN-TÊTE ---
\begin{center}
    {\Large \textbf{Loïc Aron MBASSI EWOLO}} \\ \vspace{4pt}
    \textbf{\large Élève-Ingénieur : Électronique, Microcontrôleurs \& Instrumentation} \\
    \textit{\small Disponible dès \textbf{Avril 2026} pour 4 à 6 mois} \\ \vspace{4pt}
    \small
    \faMapMarker\ Cergy, France (Mobile Île-de-France) \quad
    \faPhone\ (+33) 7 59 52 33 98 \quad
    \faEnvelope\ \href{mailto:loicaron.mbassiewolo@ensea.fr}{loicaron.mbassiewolo@ensea.fr} \\ \vspace{2pt}
    \faLinkedin\ \href{https://www.linkedin.com/in/loic-aron-mbassi-ewolo-3942a9246/}{linkedin} \quad
    \faGithub\ \href{https://github.com/Nameless0l}{github} \quad
    \faGlobe\ \href{https://mbassiloic.tech}{mbassiloic.tech}
\end{center}

\vspace{-6pt}

% --- PROFIL ---
\noindent\small{Élève-ingénieur en double diplôme à l'\textbf{ENSEA}, spécialisé en \textbf{électronique}, \textbf{systèmes embarqués} et traitement du signal. Expérience en conception de prototypes à base de \textbf{microcontrôleurs} (STM32, Arduino), programmation \textbf{Python/C} et instrumentation. Je souhaite contribuer au \textbf{CEA DAM} pour le développement du système de compteur d'impulsions de l'installation ELSA.}

% --- FORMATION ---
\section{Formation}
\begin{itemize}[leftmargin=*, label={.}, nosep]
    \item \textbf{ENSEA (École Nationale Supérieure de l'Électronique et de ses Applications)} \hfill \textit{Cergy, France | 2025 -- Présent} \\
    \small{Ingénieur 2A, Double Diplôme - \textbf{Électronique}, Traitement du Signal et Systèmes Embarqués.} \\
    \textit{\small{Cours : Électronique Analogique \& Numérique, VHDL/FPGA, Microcontrôleurs, Instrumentation, Signal.}}
    \vspace{2pt}
    \item \textbf{École Nationale Supérieure Polytechnique de Yaoundé} \hfill \textit{Cameroun | 2021 -- 2025} \\
    \small{Diplôme d'Ingénieur de Conception (Master 2) : Génie Logiciel, Mathématiques Appliquées, Algorithmique.}
    \item \textbf{Université de Yaoundé I} \hfill \textit{Cameroun | 2020 -- 2022} \\
    \small{Licence Informatique \& Mathématiques : Algorithmique C, Analyse, Algèbre Linéaire.}
\end{itemize}

% --- COMPÉTENCES TECHNIQUES ---
\section{Compétences Techniques}
\begin{itemize}[leftmargin=*, nosep]
    \item \small{\textbf{Électronique :} Analogique \& Numérique, Conception de circuits, Soudure, Capteurs (IMU, Flex).}
    \item \small{\textbf{Microcontrôleurs :} \textbf{STM32}, Arduino, Raspberry Pi, Protocoles (I2C, SPI, UART).}
    \item \small{\textbf{Programmation :} \textbf{Python} (Avancé), \textbf{C/C++}, VHDL, MATLAB/Simulink.}
    \item \small{\textbf{Outils :} Git, Linux, Oscilloscope, Multimètre, Documentation technique, Cahier des charges.}
\end{itemize}

% --- EXPÉRIENCES ---
\section{Expériences Significatives}
\begin{itemize}[leftmargin=*, label={}]

    \resumeSubheading
      {Stage Recherche - Systèmes Embarqués \& IoT}{Oct. 2023 -- Janv. 2024}
      {Laboratoire IOTA (IoT Lab) - École Polytechnique}{\textbf{Yaoundé, Cameroun}}
      \begin{itemize}[leftmargin=0.4cm, nosep, label=\tiny$\bullet$]
        \item \small{Développement de prototypes à base de \textbf{microcontrôleurs} (Arduino, Raspberry Pi).}
        \item \small{Optimisation d'algorithmes pour contraintes \textbf{temps réel} et ressources limitées.}
        \item \small{Déploiement de modèles ML embarqués avec TensorFlow Lite.}
      \end{itemize}
    \vspace{2pt}

    \resumeSubheading
      {Assistant de Recherche - Intelligence Artificielle}{Juin 2025 -- Sept. 2025}
      {MILA (Institut Québécois d'IA) - Collaboration à distance}{\textbf{Québec, Canada}}
      \begin{itemize}[leftmargin=0.4cm, nosep, label=\tiny$\bullet$]
        \item \small{Recherche sur la généralisation des réseaux de neurones sous \textbf{PyTorch}.}
        \item \small{Implémentation d'expériences reproductibles et analyse mathématique rigoureuse.}
      \end{itemize}
    \vspace{2pt}

    \resumeSubheading
      {Développeur Backend (Mission Freelance)}{Oct. 2024 -- Janv. 2025}
      {CSC Fintech - Plateforme Bancaire}{Yaoundé, Cameroun}
      \begin{itemize}[leftmargin=0.4cm, nosep, label=\tiny$\bullet$]
        \item \small{Conception d'une API critique avec gestion des accès concurrents.}
        \item \small{Rédaction de documentation technique et \textbf{cahier des charges}.}
      \end{itemize}
\end{itemize}

% --- PROJETS ---
\section{Projets Académiques \& Techniques}
\begin{itemize}[leftmargin=*, label={}]

\item \textbf{Harmony Gloves - Prototype Électronique Instrumenté} \hfill \textit{Laboratoire IOTA}
\vspace{-2pt}
    \begin{itemize}[leftmargin=0.4cm, nosep, label=\tiny$\bullet$]
        \item \small{Conception d'un circuit électronique avec \textbf{capteurs IMU et Flex sensors}.}
        \item \small{\textbf{Soudure} de composants et intégration hardware/software sur \textbf{STM32/Arduino}.}
        \item \small{Acquisition et traitement de signaux en temps réel pour classification.}
    \end{itemize}
    \vspace{3pt}

\item \textbf{Projet Fault Detection \& Control (Drone)} \hfill \textit{Projet d'option ENSEA}
\vspace{-2pt}
    \begin{itemize}[leftmargin=0.4cm, nosep, label=\tiny$\bullet$]
        \item \small{Modélisation dynamique et simulation sous \textbf{MATLAB/Simulink}.}
        \item \small{Conception de schémas de détection de défauts (observateurs, résidus).}
        \item \small{Stratégie de contrôle tolérant aux pannes avec reconfiguration.}
    \end{itemize}
    \vspace{3pt}

\item \textbf{Upgrade Jetson - Robotique Autonome (Challenge CoHoMa Défense)} \hfill \textit{En cours}
\vspace{-2pt}
    \begin{itemize}[leftmargin=0.4cm, nosep, label=\tiny$\bullet$]
        \item \small{Intégration de capteurs (Lidar 3D, caméras) sur plateforme \textbf{Nvidia Jetson}.}
        \item \small{Optimisation d'algorithmes de navigation autonome (SLAM).}
    \end{itemize}
    \vspace{3pt}

\item \textbf{Projet Taxi - Simulation Temps Réel en C} \hfill \textit{Projet Académique}
\vspace{-2pt}
    \begin{itemize}[leftmargin=0.4cm, nosep, label=\tiny$\bullet$]
        \item \small{Développement d'un simulateur multiprocessus avec gestion des \textbf{impulsions/signaux}.}
        \item \small{Synchronisation par signaux UNIX, mémoire partagée (SHM), ordonnanceur FIFO.}
        \item \small{Visualisation 2D temps réel avec OpenGL.}
    \end{itemize}
    \vspace{3pt}

\item \textbf{Projet UROP - Analyse de Signaux ECG} \hfill \textit{Collaboration Google DeepMind}
\vspace{-2pt}
    \begin{itemize}[leftmargin=0.4cm, nosep, label=\tiny$\bullet$]
        \item \small{Traitement de signaux médicaux : \textbf{filtrage numérique}, extraction de features.}
        \item \small{Modèle de classification pour détection d'anomalies.}
    \end{itemize}

\end{itemize}

% --- DISTINCTIONS ---
\section{Distinctions, Certifications \& Langues}
\begin{itemize}[leftmargin=*, nosep]
    \item \textbf{Hackathons :} \textbf{1\textsuperscript{ère} Place} IndabaX Africa, \textbf{1\textsuperscript{ère} Place} CITS National, \textbf{1\textsuperscript{ère} Place} Mountain Hub.
    \item \textbf{Leadership :} Chef Cellule IA \& Projets (Polytechnique) - Animation workshops, gestion d'équipes.
    \item \textbf{Certifications :} Supervised ML (DeepLearning.AI/Stanford), Python for Data Science (IBM).
    \item \textbf{Langues :} Français (Natif), Anglais (Professionnel/Technique).
    \item \textbf{Académique :} Note Baccalauréat Physique : 19.25/20.
\end{itemize}

\end{document}
